\chapter{Optimization II: memory, caching and data layout}\label{ch:memory}

Most of WRF on a GPU is limited by memory bandwidth (\cref{ch:perf}). A kernel can therefore become faster in two
ways only: by moving fewer bytes, or by moving them more efficiently. This chapter covers both: how WRF's array
layout meets the GPU's memory system, how to keep data in the caches between uses, and how to use the special memory
paths a GPU offers. None of it changes an arithmetic statement, so all of it is class R0 or R1 (\cref{tab:bitclass}).

\section{Coalescing and WRF's layout}

A warp's load is served in 32-byte sectors. It is \emph{coalesced} when the 32 threads read 32 consecutive 4-byte
values, which costs four sectors (\cref{ch:gpuarch}). The profiler reports the average number of sectors per request;
four is ideal, 32 is the worst case.

\paragraph{Fields.} WRF stores a 3D field as \code{a(ims:ime, kms:kme, jms:jme)}: $i$ is contiguous, then $k$, then
$j$. A pointwise kernel collapses the source's loops \code{j}, \code{k}, \code{i} into one iteration space, and
consecutive threads get consecutive $i$. The loads are coalesced, apart from the warps that straddle the end of one
row and the start of the next, which need one or two extra sectors. The 4D arrays (moisture species, scalars) have the
species index last, so each species is a 3D field of its own, with the same layout.

\paragraph{Columns.} A column kernel gives each thread an $(i,j)$ column and loops over $k$ inside the thread. At each
level $k$ the 32 threads of a warp read 32 consecutive $i$: coalesced again. This is the reason WRF's layout suits
GPUs for both kinds of kernel. The layout that seems natural for column physics, with $k$ contiguous, would make each
thread's column contiguous, but the threads of a warp would then read addresses a whole column apart, and every load
would be uncoalesced.

\paragraph{Private arrays.} The private column arrays of a column kernel live in local memory. The compiler
interleaves local memory across threads: element $m$ of thread $t$ lies next to element $m$ of thread $t+1$. When all
threads of a warp access the same element (the same level $k$), the access is coalesced.

\paragraph{Exceptions.} A row-parallel kernel (Template R) gives each thread a row $j$, and loops over $i$ inside. The
threads of a warp then read addresses one whole $(i,k)$ plane apart: 32 sectors per request. The ozone interpolation
is such a kernel. It runs only on d01 and only at radiation steps, so it costs little; if the profiler finds it
expensive, the remedy is to show (test T-OZN) that a per-column search finds the same bracket as the source's shared
search, and to give each column its own thread.

\section{Reuse through the caches}

A value read twice from device memory costs twice; a value read once and found in a cache the second time costs once.
Whether the second read hits depends on how much data passes through the cache between the two reads, the
\emph{reuse distance}:
\begin{itemize}
  \item \textbf{Within a kernel.} A stencil reads each value several times, from neighbouring threads. Neighbours in
    $i$ are in the same warp and share sectors. Neighbours in $j$ or $k$ are in other warps or blocks, and their reads
    hit only if the data is still in L1 or L2. The 5th-order flux in $y$ reads six rows of the field. If those rows
    are not in a cache, the kernel moves six times the minimal bytes. Tiling (\cref{ch:tiling}) makes the reuse
    explicit.
  \item \textbf{Between kernels.} A kernel that reads what the previous kernel wrote finds it in L2 if the data fits.
    A d02 3D field of the full case (162~MB) never fits; a dev-case field (7.9~MB) does. This is a second reason why
    dev-case timings do not carry over to the full case, and a reason why fusion (\cref{ch:fusion}) saves more on the
    full case: there, the intermediate array really goes to device memory and back.
\end{itemize}

\subsection{L2 residency control}

The A100 and H100 let a program set aside part of the L2 cache for \emph{persisting} accesses: an address window,
attached to a CUDA stream, whose lines the cache prefers to keep (an access policy window). Candidates are arrays that
many kernels read again and again and that fit:
\begin{itemize}
  \item the absorption tables of RRTMG (about 0.75~MB), read by every column at every radiation step;
  \item the lookup tables of the surface layer and the Noah parameters;
  \item the fire grid's level-set function, read by every kernel of the fire step. One full-case fire field
    (43~MB) does not fit the A100's 40~MB L2 at all, and only just fits the H100's 50~MB.
\end{itemize}
Phase S6-09 measures the effect on both GPUs. The change is in class R0: the cache policy changes where a value is
read from, never its bits.

\subsection{The cache directive}

OpenACC's \code{cache} directive asks the compiler to keep a range of an array in the fastest memory available
inside a loop, for example the six values a 5th-order stencil reads along one direction. It is a hint: how much the
compiler does with it varies, and only the measured sectors per request and DRAM bytes show whether it helped.
Where it does not, the same staging is written by hand in CUDA Fortran (\cref{ch:tiling}).

\section{Special memory paths}

\paragraph{Read-only path.} Loads of data that the kernel never writes can go through the read-only (texture) data
path, which has its own cache and handles scattered access well. The compiler uses it when it can prove that the
array is not written in the kernel. Listing arrays as read-only in CUDA Fortran (\code{intent(in)} with the
\code{device} attribute) makes that proof easy.

\paragraph{Constant memory.} A small read-only area of 64~KB, cached separately, which serves all threads of a warp in
one transaction when they read the \emph{same} address. It suits constants that every thread reads (WSM6's
coefficients, the fire flags). It does not suit tables indexed by data, such as the stability-function tables of the
surface layer, which are indexed by the stability parameter $z/L$ of each point: when the threads of a warp read
different addresses, constant memory serializes them. Those tables are better left in device memory, where L1
serves them.

\paragraph{Shared memory.} The memory a block manages itself: the basis of tiling (\cref{ch:tiling}).

\section{Transfers to the host}

Copies between host and device go fastest from \emph{pinned} host memory, which the operating system may not move.
From ordinary (pageable) memory, the driver copies through a pinned staging buffer, which halves the bandwidth and
prevents overlap with computation. The history output of d02, written every 15 simulated minutes, is the largest
regular copy to the host. Optimization O8 copies it into pinned buffers asynchronously, so that the GPU continues with
the next steps while the host writes the file. The values written are the same; only the timing changes.

\section{Worked examples}

The profiler's analyzers A2 (caching) and A4 (coalescing) produce the candidate lists. Phases S6-09 and S6-10 take them
in order:
\begin{enumerate}
  \item \textbf{Advection.} The $y$-flux kernels read six rows of the advected field per face. The measured DRAM bytes
    against the minimal bytes show how much the caches already catch; the \code{cache} directive and then the tile of
    \cref{ch:tiling} are tried in turn.
  \item \textbf{Fire WENO.} \code{tend_ls} reads a $\pm 3$ neighbourhood of the level-set function in both directions.
    On the full case, L2 persistence of the level-set field is tried on the H100, where it fits.
  \item \textbf{Coalescing.} Kernels with more than 4.5 sectors per request are listed and their access patterns
    examined; each fix is one kernel per commit.
\end{enumerate}
The numbers before and after each step are \tbr.

\begin{keypoint}
WRF's layout, with $i$ contiguous, is already the right one for a GPU: it coalesces both the pointwise stencils and
the column kernels. Most memory optimizations of the port are therefore about reuse (moving each byte fewer times),
not about layout.
\end{keypoint}
